\begin{table}[H]
\centering
\caption{Summary Statistics}
\label{tab:summary}
\begin{threeparttable}
\begin{tabular}{lrrrrr}
\toprule
 & N & Mean & Std.\ Dev. & Min & Max \\
\midrule
\multicolumn{6}{l}{\textit{Panel A: Full Sample}} \\
CPI Index & 19,493 & 113.4 & 18.8 & 52.2 & 341.3 \\
log(CPI) & 19,493 & 4.719 & 0.157 & 3.955 & 5.833 \\
\midrule
\multicolumn{6}{l}{\textit{Panel B: By Tax Treatment Group}} \\
Group A: GST-taxed, SST-covered & 3,860 & 112.8 & 16.0 & 71.6 & 202.1 \\
Group B: GST-taxed, no SST & 7,913 & 108.5 & 18.9 & 52.2 & 341.3 \\
Group C: Zero-rated/Exempt & 7,720 & 119.1 & 18.4 & 57.1 & 191.0 \\
\midrule
\multicolumn{6}{l}{\textit{Panel C: By Period}} \\
Pre-GST (2010--2015) & 6,363 & 103.8 & 6.3 & 68.2 & 142.2 \\
GST era (2015--2018) & 3,838 & 113.3 & 13.8 & 62.6 & 179.3 \\
Tax holiday (Jun--Aug 2018) & 303 & 113.6 & 16.1 & 59.2 & 179.3 \\
SST era (Sep 2018+) & 8,989 & 120.2 & 23.0 & 52.2 & 341.3 \\
\bottomrule
\end{tabular}
\begin{tablenotes}[flushleft]
\small
\item Notes: 19,493 product-month observations across 101 COICOP 4-digit product classes and 193 months (January 2010--January 2026). CPI indices from OpenDOSM with base year 2010 = 100. Group A products were standard-rated at 6\% under GST (April 2015--May 2018) and are covered by SST (September 2018 onward). Group B products were standard-rated under GST but are not covered by SST. Group C products were zero-rated or exempt under GST and serve as controls.
\end{tablenotes}
\end{threeparttable}
\end{table}

